\subsection{Derived \texorpdfstring{$48=4\cdot12$}{48=4·12} identity and the \texorpdfstring{$2$-$3$-$5$}{2-3-5} arithmetic}

\begin{proposition}[Derived \texorpdfstring{$48=4\cdot12$}{48=4·12} identity on the retained branch]
\label{prop:rank-lock}
On the retained branch, once the family multiplicity and carrier ranks have been fixed, one
recovers
\[
\Oadm
=
N_{\mathrm{fam}}\,\dim S^+
=
(g-1)\dim\mathfrak g_{\mathrm{car}}
=
(g-1)(b^2+s^2-1)
\]
with
\[
N_{\mathrm{fam}}=3,
\qquad
\dim S^+=16,
\qquad
(b,s)=(3,2),
\qquad
g=5.
\]
Equivalently, the former rank-lock formula is a derived consistency identity evaluated at the
closed branch point:
\begin{equation}
\label{eq:rank-lock}
3\cdot 2^{g-1}=(g-1)(b^2+s^2-1).
\end{equation}
\end{proposition}

\begin{proof}
\Cref{thm:harmonic-family-modes} gives the physical family multiplicity
\[
N_{\mathrm{fam}}=\dim\mathcal H_F=3.
\]
The carrier closure is completed later by \TFPTxref{cor:bosonic-rank-two,thm:yukawa-forces-32}, which
fix
\[
(b,s)=(3,2),
\qquad
g=5,
\qquad
\dim S^+=16.
\]
Therefore
\[
\Oadm=N_{\mathrm{fam}}\dim S^+=3\cdot16=48.
\]
For the same minimal branch,
\[
\dim\mathfrak g_{\mathrm{car}}
=
\dim\mathfrak{su}(3)+\dim\mathfrak{su}(2)+\dim\mathfrak u(1)
=
12.
\]
Therefore
\[
\Oadm=48=(g-1)\dim\mathfrak g_{\mathrm{car}}=4\cdot 12,
\]
which is exactly the displayed derived identity written in branch variables. The former occupancy
balance is thus a consistency check rather than an input assumption.
\end{proof}

\begin{remark}[The $2$-$3$-$5$ arithmetic]
The hard carrier packet rests on exactly three discrete numbers: $2$ (the double cover), $3$ (the color rank and family count), and $5$ (the carrier rank). From these three atoms one recovers
\[
\gamma=\frac{5}{6}=\frac{5}{2\cdot 3},
\qquad
B=\frac{3}{2},
\qquad
B\gamma=\frac{5}{4}=\frac{5}{2^2},
\]
\[
\Oadm=3\cdot 2^4=48,
\qquad
\Dstart=5\cdot 12=60,
\qquad
\frac{\Dstart}{\Oadm}=\frac{5}{4}=B\gamma,
\]
and the defect coefficients $\deltatop=48\,\cthree^4$ and $\delta_2=\frac54\,\deltatop^2=2880\,\cthree^8$ with $2880=60\cdot 48$.
The persistent appearance of $5/4$ as carrier asymmetry $\times$ norm, second defect stage, and ladder-to-occupancy ratio is therefore not three coincidences but one rational shadow of the $2$-$3$-$5$ triad.
\end{remark}

\begin{lemma}[Universal \texorpdfstring{$5/4$}{5/4} compression quotient on the minimal branch]
\label{lem:universal-5-4-compression}
Define
\[
\rho_{5/4}:=B\gamma.
\]
On the canonical minimal branch,
\[
\rho_{5/4}
=
B\gamma
=
\frac{\delta_2}{\deltatop^2}
=
\frac{\Dstart}{\Oadm}
=
\frac{\gcar}{\gcar-1}
=
\frac54.
\]
Thus $5/4$ is the universal compression quotient of the minimal branch rather than a recurring numerical motif.
\end{lemma}

\begin{proof}
The identity $B\gamma=\frac54$ is \Cref{cor:minimal-carrier-compression}. The defect relation gives
\[
\frac{\delta_2}{\deltatop^2}=\frac54.
\]
On the canonical branch,
\[
\frac{\Dstart}{\Oadm}=\frac{60}{48}=\frac54,
\]
and with $\gcar=5$ one also has
\[
\frac{\gcar}{\gcar-1}=\frac{5}{4}.
\]
Hence all displayed ratios agree on the same exact value.
\end{proof}

\begin{remark}[Interpretive status of the quotient]
The quotient itself is exact. Any holographic or spacetime-dimensional reading of the denominator $\gcar-1=4$ remains interpretive in the present manuscript rather than theorem-level.
\end{remark}

\begin{corollary}[Rank-5 master compression on the retained branch]
\label{cor:rank-5-master-compression}
On the retained branch one has $\gcar=5$. Consequently
\[
N_{\mathrm{fam}}=\frac{\gcar+1}{2}=3,
\qquad
\Oadm=(\gcar+1)\,2^{\gcar-2}=48,
\]
\[
\Oadm\gamma=\gcar\,2^{\gcar-2}=40,
\qquad
\bone=\frac{\gcar\,2^{\gcar-2}+1}{10}=\frac{41}{10},
\qquad
10\,\bone=41.
\]
Thus the closed family count, admissible occupancy, weighted occupancy, and abelian index all
compress to the single carrier rank.
\end{corollary}

\begin{proof}
By \TFPTxref{cor:bosonic-rank-two,thm:yukawa-forces-32}, the retained carrier has
\[
\gcar=5.
\]
By \Cref{thm:harmonic-family-modes}, the family block fixes $N_{\mathrm{fam}}=3$, hence
\[
\Oadm=N_{\mathrm{fam}}\,2^{\gcar-1}=3\cdot 2^4=48.
\]
The carrier norm corollary gives $\gamma=5/6$, so
\[
\Oadm\gamma=48\cdot\frac56=40.
\]
Finally the closed Higgs branch has $N_\Phi=1$, and \TFPTeqref{eq:41-chain-main} yields
\[
10\,\bone=\Oadm\gamma+N_\Phi=40+1=41.
\]
\end{proof}

\begin{corollary}[Pythagorean compression of the abelian index on the retained branch]
\label{cor:pythagorean-integrity-index}
On the retained branch,
\[
\mathcal I_{41}
=
10\,\bone
=
2\gcar(\gcar-1)+1
=
\gcar^2+(\gcar-1)^2.
\]
In particular, since $\gcar=5$ on the retained branch, one gets
\[
\mathcal I_{41}=5^2+4^2=41.
\]
\end{corollary}

\begin{proof}
By \TFPTxref{cor:rank-5-master-compression}, one has $\gcar=5$ and
\[
10\,\bone=\gcar\,2^{\gcar-2}+1=5\cdot 2^3+1=41.
\]
\[
2\gcar(\gcar-1)+1=2\cdot 5\cdot 4+1=41.
\]
Since on the retained branch
\[
2\gcar(\gcar-1)+1=\gcar^2+(\gcar-1)^2,
\]
the claim follows.
\end{proof}

\begin{remark}[The number $4$ as a cross-cutting echo]
In addition to the $2$-$3$-$5$ triad, the integer $4$ appears at four structurally distinct levels:
\[
F_{\mathrm{patch}}=4,
\qquad
\gcar-1=4,
\qquad
\frac{\Oadm}{\dim\mathfrak g_{\SM}}=\frac{48}{12}=4,
\qquad
\frac{|R(E_8)|}{\Dstart}=\frac{240}{60}=4.
\]
That is, four corner patches, carrier rank minus one, matter-to-gauge compression factor, and root-to-start quotient all evaluate to the same integer.
\end{remark}

